\section{Introduction}
\label{sec:introduction}


% 1
Coding agents are increasingly deployed on long-horizon software-engineering tasks, with measurable progress on issue resolution over real-world code repositories~\citep{jimenezSWEbenchCanLanguage2023,yangSWEbenchMultimodalAI2024,dengSWEBenchProCan2025} and multi-step terminal workflows~\citep{merrillTerminalBenchBenchmarkingAgents2026a}.
In practice, such progress relies not only on the underlying language model, but equally on the surrounding engineering components: the system prompt that shapes work style, the tools that expose the file system and shell, and the middleware that controls context, execution, and recovery.
This collection of model-external, editable components is collectively referred to as the agent's \emph{harness}~\citep{rajasekaranHarnessDesignLongrunning2026,lopopoloHarnessEngineeringLeveraging2026,wangOpenHandsOpenPlatform2025,yangSWEagentAgentComputerInterfaces2024,steinbergerOpenClawPersonalAI2026,HermesAgentAgent}.

% 2

Harness design materially shifts task completion on long-horizon coding benchmarks, even with the base model held fixed~\citep{trivedyImprovingDeepAgents2026,wangOpenHandsOpenPlatform2025}, making harness engineering a first-class lever for improving coding agents.
Moreover, the optimal harness is model-specific: a harness tuned for one base model often underperforms on another and must be re-adapted as the base model changes.
In current practice, this adaptation is performed manually—developers inspect trajectories, identify recurring failure patterns, and hand-craft edits across prompts, tools, middleware, and skills.
Yet as base models advance rapidly~\citep{xiaomimimoteamMiMoV25Pro2026,teamQwen36PlusRealWorld2026,yangQwen3TechnicalReport2025a,DeepSeek_V4pdfDeepseekaiDeepSeekV4Pro2026,kimiteamKimiK26Tech2026,teamKimiK25Visual2026}, this manual loop struggles to keep pace, creating a widening gap between model capability and the harness needed to realize it~\citep{steinbergerOpenClawPersonalAI2026}.

% 3 

An intuitive direction is to automate this loop with an \textit{evolution agent} that optimizes harness components based on experience~\citep{agrawalGEPAReflectivePrompt2025a,zhangAgenticContextEngineering2025a,caiTrainingFreeGroupRelative2025}.
However, few existing approaches jointly evolve the full set of editable components~\citep{leeMetaHarnessEndtoEndOptimization2026}; most focus on a single component, typically the prompt~\citep{shinnReflexionLanguageAgents2023c,zhaoExpeLLLMAgents2024,madaanSelfRefineIterativeRefinement2023c}, skills~\citep{maSkillClawLetSkills2026b,xiaSkillRLEvolvingAgents2026c}, or an in-context playbook~\citep{zhangAgenticContextEngineering2025a}.
Jointly evolving multiple components end-to-end faces two structural obstacles: long, unstructured trajectories yield little actionable signal, and tightly coupled harness frameworks make edits beyond the prompt error-prone.
This leaves the central question of agent-driven harness evolution open:
\textit{How can an evolution agent \textbf{jointly} and \textbf{stably} evolve all editable components of a coding agent's harness?}

% 4

Our central insight is that this question is bottlenecked by \emph{observability}, not by agent capability: once the evolution agent receives structured context over a clear action space, it can reliably converge on better harness designs~\citep{suttonBitterLesson2019,zunicBitterLessonAgent2026}.
We implement this in \textbf{A}gentic \textbf{H}arness \textbf{E}ngineering 
(\textbf{AHE}, Figure~\ref{fig:method}), a closed loop driven by three observability pillars:
\ding{182}~\emph{component observability} via a decoupled harness that exposes seven editable component types as files, so each failure pattern maps cleanly to a single component class;
\ding{183}~\emph{experience observability} via a layered, drill-down evidence corpus distilled from millions of raw trajectory tokens, so the evolver consumes structured root causes rather than raw logs; and
\ding{184}~\emph{decision observability} via a change manifest that pairs every edit with a self-declared prediction, later verified against the next round's task-level outcomes, so each edit becomes a falsifiable contract and ineffective ones are reverted at file granularity.

% 5

We empirically validate AHE on Terminal-Bench 2\citep{merrillTerminalBenchBenchmarkingAgents2026a}: ten iterations lift pass@1 from 69.7\% to 77.0\%, surpassing the human-designed Codex~\citep{openaiCodexCLI2025} and the self-evolving baselines ACE~\citep{zhangAgenticContextEngineering2025a} and Training-Free GRPO~\citep{caiTrainingFreeGroupRelative2025}.
% Without further evolution, the frozen harness transfers to SWE-bench-verified~\citep{jimenezSWEbenchCanLanguage2023} and to three alternate base-model families, with the largest gains on bases further from saturation, suggesting that AHE encodes coordination patterns that less-saturated models lean on more heavily.
Without further evolution, the frozen harness transfers to SWE-bench-verified~\citep{jimenezSWEbenchCanLanguage2023}, and across three alternate base-model families it yields consistent pass@1 gains of $+5.1$ to $+10.1$ pp, with the largest on bases further from saturation, suggesting that AHE encodes coordination patterns that less-saturated models lean on more heavily.
A component ablation pinpoints where this gain lives: tools, middleware, and long-term memory each carry the improvement on their own, while the system prompt alone regresses, indicating that factual harness structure transfers across tasks and models whereas prose-level strategy does not.

This paper makes three contributions:
\begin{itemize}
    \item We formulate \emph{agent-driven harness evolution} for coding agents and propose \textbf{AHE}, which identifies \emph{observability across components, trajectories, and decisions} as the design pivot and turns every harness edit into a falsifiable, file-level contract through three observability pillars: a decoupled component substrate, a layered trajectory-distillation pipeline, and a change manifest whose self-declared predictions are verified by next-round task deltas.
    \item We empirically show that AHE lifts pass@1 on Terminal-Bench~2 from 69.7\% to 77.0\%, surpasses human-designed and automated baselines, and produces a frozen harness that transfers across benchmarks and base-model families.
    \item Our analysis reveals two limits of agent-driven evolution: harness components interact non-additively, so stacking effective edits caps the aggregate gain; and the loop's self-attribution is reliable for fixes but blind to regressions, pinpointing regression foresight as the clearest direction for future self-evolution loops.
\end{itemize}

